\chapter{指数与对数.}

实数乘法群 \(R_{+}^{*}\)（\textgreater0）的历史，与一个数的幂（\textgreater0）这一观念的发展以及用来表示它们的记法密切相关。由同一个数的逐次幂组成的“几何级数”这一概念，可以追溯到埃及人和巴比伦人；希腊数学家对它也很熟悉，而且在欧几里得（《几何原本》第九卷，命题 11）那里就已经可以找到一个一般命题，它与规则 \(a^{m}a^{n}=a^{m+n}\)（指数为 \textgreater0 的整数）等价。在中世纪，法国数学家 N. 奥雷姆（十四世纪）重新发现了这条规则；分数指数（\textgreater0）的概念也首次出现在他那里，其记法已经与我们的相近，而且还伴有一些以一般形式陈述的关于它们的运算法则（例如我们现在写作 \((ab)^{1/n}=a^{1/n}b^{1/n},(a^{m})^{p/q}=(a^{mp})^{1/q}\) 的那两条法则）{[}74{]}。但奥雷姆的思想远远超前于他那个时代的数学，以致未能对其同时代人发生影响，他的论著很快就被遗忘了。一个世纪之后，N. 许凯再次陈述了欧几里得的法则；他还在自己的方程中为未知量的幂引入了一种指数记法，并且毫不犹豫地使用了指数 0 和负整数指数（\textless0）。¹这一次（尽管许凯的著作一直停留在手稿状态，似乎也未曾广泛流传），指数的“算术级数”与幂的“几何级数”之间的同构这一观念再也没有从人们的视野中消失；它被施蒂费尔推广到负指数和分数指数（{[}298{]}，fol.~35、249-250），最终导致了 对数的定义和最初数表的编制，后者由苏格兰人 J. 纳皮尔于 1614-1620 年间、瑞士人 J. 比尔吉（其著作迟至 1620 年才问世，尽管其构想可追溯到十七世纪的最初几年）各自独立地完成。在比尔吉那里，通过在处理数表时使用插值，R 与 \(R_{+}^{*}\) 之间所建立的同构的连续性是被隐含地假定的；相反，它被明确地表述于

纳皮尔的定义之中（至少在当时通行的、还相当模糊的连续性概念所允许的程度上是明确的）。\(^{2}\)

我们在这里无需详述对数在数值计算中所提供的服务；从理论观点看，它们的重要性尤其始于无穷小计算的萌芽时期，即 \(\log(1+x)\) 与 \(e^{x}\) 的级数展开以及这些函数的微分性质被发现之时（见第 172 页及以下）。至于指数函数和对数函数的定义，直到十九世纪中叶，人们都仅限于直观地假定：有可能把对一切有理数 x 定义的函数 \(a^{x}\) 连续地延拓到实数集上；而只是在实数概念被最终精确化并从有理数概念导出之后，人们才想到要对这一延拓给出严格的证成。

